% Generated by output/ablation_revision_20261011/build_appendix.py.
% Frozen values SHA256: 5a9ce882749c622d5ebe58a640b36c94d98d7d38117f413b2aad8229282a0287
\begin{ablationadded}
\subsection{Dynamic Core Replacement Across Backbones}
\label{app:new_core_ablation}

We extend paired core replay to Llama 3.1-8B-Instruct, Qwen3.5-9B, and Kimi-VL-A3B, with one Full run per model--case. Each model uses the same 80 injected variants: four templates applied to 20 source SCCs across CUAD, SEC, BGB, and Debian. Each run uses 60 rollouts, four-node windows, internal concurrency four, and temperature 0.4. These 240 trajectories form the ablation cohort, separate from the historical main-experiment results. Tables in this and the next three subsections report percentages averaged over 80 cases per model; the 20 source SCCs are the independent sampling units.

To isolate replacement, both policies receive the same ordered evidence proposals. Dynamic allows earlier admissions to be replaced; locked retains admitted nodes without eviction. For each case, locked capacity is fixed to the final Dynamic output size $k$, and both return exactly $k$ nodes. This paired replay isolates replacement at fixed evidence and retained-core size. Root@3 uses the shared trajectory ranking and is identical across selectors.

\begin{table}[!htbp]
\ablationcolor
\centering
\small
\renewcommand{\arraystretch}{1.10}
\setlength{\tabcolsep}{3pt}
\begin{tabular}{lrrrr}
\toprule
\bfseries Policy & \bfseries Root@3 & \bfseries Any & \bfseries All & \bfseries Comp. \\
\midrule
\multicolumn{5}{l}{\bfseries Llama 3.1-8B-Instruct} \\
Dynamic (Full) & \tblbest{62.50} & \tblbest{93.75} & \tblbest{62.50} & 51.26 \\
Locked & \tblbest{62.50} & \tblsecond{85.00} & \tblsecond{35.00} & 51.26 \\
\midrule
\multicolumn{5}{l}{\bfseries Qwen3.5-9B} \\
Dynamic (Full) & \tblbest{58.75} & \tblbest{88.75} & \tblbest{57.50} & 66.55 \\
Locked & \tblbest{58.75} & \tblsecond{70.00} & \tblsecond{33.75} & 66.55 \\
\midrule
\multicolumn{5}{l}{\bfseries Kimi-VL-A3B} \\
Dynamic (Full) & \tblbest{7.50} & \tblbest{93.75} & \tblbest{35.00} & 51.93 \\
Locked & \tblbest{7.50} & \tblsecond{80.00} & \tblsecond{30.00} & 51.93 \\
\bottomrule
\end{tabular}
\caption{Dynamic versus locked core replay at the default capacity ratio $r=0.45$. Any and All denote Risk-any and Risk-all; Comp. is $100(1-|H|/|S|)$. Evidence and actual output size are matched per case. Bold/underline mark the best/second-best endpoint rate within each model, including ties; compression is descriptive.}
\label{tab:new_core_replacement}
\end{table}

Dynamic replacement increases Risk-all by 27.50, 23.75, and 5.00 percentage points for Llama, Qwen, and Kimi, respectively (Table~\ref{tab:new_core_replacement}). Risk-any also improves on all three backbones at identical compression. The comparison supports revising the retained core as evidence accumulates, with the largest gains on Llama and Qwen.

\subsection{Core Capacity and Retention}
\label{app:new_core_capacity}

We replay the same Full evidence at four capacity ratios, $r\in\{0.25,0.35,0.45,0.60\}$, with cap $\min\{|S|,\max(4,\lceil r|S|\rceil)\}$. At each ratio, we recompute the production selector from the recorded evidence. Locked admission matches Dynamic's actual final $k$ case by case. The cap is an upper bound; observed compression is calculated from the actual retained-core size.

\begin{table}[!htbp]
\ablationcolor
\centering
\small
\renewcommand{\arraystretch}{1.10}
\setlength{\tabcolsep}{3pt}
\begin{tabular}{lrrrr}
\toprule
\bfseries $r$ & \bfseries Dynamic & \bfseries Locked & \bfseries $\Delta$ & \bfseries Comp. \\
\midrule
\multicolumn{5}{l}{\bfseries Llama 3.1-8B-Instruct} \\
0.25 & \tblbest{25.00} & 13.75 & $+11.25$ & 68.30 \\
0.35 & \tblbest{38.75} & 16.25 & $+22.50$ & 61.76 \\
0.45 & \tblbest{62.50} & 35.00 & $+27.50$ & 51.26 \\
0.60 & \tblbest{75.00} & 60.00 & $+15.00$ & 36.93 \\
\midrule
\multicolumn{5}{l}{\bfseries Qwen3.5-9B} \\
0.25 & \tblbest{47.50} & 22.50 & $+25.00$ & 73.49 \\
0.35 & \tblbest{48.75} & 26.25 & $+22.50$ & 70.31 \\
0.45 & \tblbest{57.50} & 33.75 & $+23.75$ & 66.55 \\
0.60 & \tblbest{57.50} & 38.75 & $+18.75$ & 64.21 \\
\midrule
\multicolumn{5}{l}{\bfseries Kimi-VL-A3B} \\
0.25 & \tblbest{17.50} & 7.50 & $+10.00$ & 68.34 \\
0.35 & \tblbest{20.00} & 13.75 & $+6.25$ & 61.91 \\
0.45 & \tblbest{35.00} & 30.00 & $+5.00$ & 51.93 \\
0.60 & \tblbest{55.00} & 47.50 & $+7.50$ & 38.67 \\
\bottomrule
\end{tabular}
\caption{Capacity sweep on the same 80 Full trajectories per model. Dynamic and Locked columns report Risk-all; $\Delta$ is their difference in percentage points. Comp. is the shared mean compression of the paired outputs. Bold marks the higher Risk-all within each pair.}
\label{tab:new_core_capacity}
\end{table}

Dynamic retains complete endpoints more often than locked admission at all 12 model--capacity settings (Table~\ref{tab:new_core_capacity}). Dynamic's Risk-all increases with capacity through $r=0.60$ for Llama and Kimi; Qwen reaches 57.50\% at both $r=0.45$ and $0.60$. The paired gains persist across the tested caps, while the absolute coverage and compression quantify the trade-off in output size.

\subsection{Selection from Shared Evidence}
\label{app:new_selector_control}

We compare the Full selector with LEA-score while holding the acquired evidence and final output size fixed. LEA-score applies the membership formula in Appendix~\ref{app:revision_lea} to the recorded Full evaluations and returns its top $k$ nodes, where $k$ is Full's actual output size for that case. Its signals are local-subgraph membership, repair-entry and conflict-endpoint counts, plus mean-risk surplus. The comparison isolates final selection on evidence acquired by SA-MCGS.

\begin{table}[!htbp]
\ablationcolor
\centering
\small
\renewcommand{\arraystretch}{1.10}
\setlength{\tabcolsep}{3pt}
\begin{tabular}{lrrrr}
\toprule
\bfseries Policy & \bfseries Root@3 & \bfseries Any & \bfseries All & \bfseries Comp. \\
\midrule
\multicolumn{5}{l}{\bfseries Llama 3.1-8B-Instruct} \\
Full & \tblbest{62.50} & \tblsecond{93.75} & \tblbest{62.50} & 51.26 \\
LEA-score & \tblbest{62.50} & \tblbest{96.25} & \tblsecond{57.50} & 51.26 \\
\midrule
\multicolumn{5}{l}{\bfseries Qwen3.5-9B} \\
Full & \tblbest{58.75} & \tblbest{88.75} & \tblbest{57.50} & 66.55 \\
LEA-score & \tblbest{58.75} & \tblsecond{86.25} & \tblsecond{52.50} & 66.55 \\
\midrule
\multicolumn{5}{l}{\bfseries Kimi-VL-A3B} \\
Full & \tblbest{7.50} & \tblbest{93.75} & \tblbest{35.00} & 51.93 \\
LEA-score & \tblbest{7.50} & \tblsecond{78.75} & \tblsecond{26.25} & 51.93 \\
\bottomrule
\end{tabular}
\caption{Full versus LEA-score at matched evidence and actual output size, using the default capacity. Endpoint-rate highlighting follows Table~\ref{tab:new_core_replacement}. Root@3 uses the shared trajectory ranking.}
\label{tab:new_selector_lea}
\end{table}

Full improves Risk-all over LEA-score by 5.00, 5.00, and 8.75 points on Llama, Qwen, and Kimi (Table~\ref{tab:new_selector_lea}).

\subsection{Evidence Accumulation Across Rollout Prefixes}
\label{app:new_prefix_analysis}

We reconstruct the native Full core after 15, 30, 45, and 60 completed rollout events from each stored 60-rollout trajectory. These four checkpoints measure evidence accumulation within each trajectory, with the selector, evidence thresholds, and nominal capacity ratio held fixed.

\begin{table}[!htbp]
\ablationcolor
\centering
\small
\renewcommand{\arraystretch}{1.10}
\setlength{\tabcolsep}{3pt}
\begin{tabular}{lrrrr}
\toprule
\bfseries Prefix & \bfseries Root@3 & \bfseries Any & \bfseries All & \bfseries Comp. \\
\midrule
\multicolumn{5}{l}{\bfseries Llama 3.1-8B-Instruct} \\
15 & 45.00 & 85.00 & 51.25 & 56.39 \\
30 & \tblbest{63.75} & \tblsecond{93.75} & \tblsecond{61.25} & 51.86 \\
45 & 57.50 & \tblbest{96.25} & 58.75 & 51.33 \\
60 & \tblsecond{62.50} & \tblsecond{93.75} & \tblbest{62.50} & 51.26 \\
\midrule
\multicolumn{5}{l}{\bfseries Qwen3.5-9B} \\
15 & 42.50 & 67.50 & 33.75 & 78.23 \\
30 & 55.00 & 76.25 & 45.00 & 73.19 \\
45 & \tblbest{60.00} & \tblsecond{86.25} & \tblsecond{51.25} & 69.02 \\
60 & \tblsecond{58.75} & \tblbest{88.75} & \tblbest{57.50} & 66.55 \\
\midrule
\multicolumn{5}{l}{\bfseries Kimi-VL-A3B} \\
15 & \tblsecond{5.00} & 73.75 & 13.75 & 63.14 \\
30 & \tblsecond{5.00} & 83.75 & 26.25 & 56.71 \\
45 & \tblbest{7.50} & \tblsecond{91.25} & \tblsecond{31.25} & 53.03 \\
60 & \tblbest{7.50} & \tblbest{93.75} & \tblbest{35.00} & 51.93 \\
\bottomrule
\end{tabular}
\caption{Native Full outputs reconstructed at four rollout prefixes. Each row averages the same 80 cases for its model. Root@3, Risk-any (Any), Risk-all (All), and compression (Comp.) are percentages. Bold/underline identify the highest/second-highest distinct endpoint rate across prefixes within each model; compression describes the evolving output size.}
\label{tab:new_native_prefixes}
\end{table}

From 15 to 60 events, native Risk-all rises from 51.25\% to 62.50\% on Llama, 33.75\% to 57.50\% on Qwen, and 13.75\% to 35.00\% on Kimi (Table~\ref{tab:new_native_prefixes}). The corresponding cores retain more nodes, quantifying the coverage--compression trade-off as search accumulates evidence. This trajectory analysis complements the paired replacement comparisons at fixed evidence and output size.

\end{ablationadded}
